\chapter{Introduction}\label{sec:introduction}

\section{Main result}\label{sec:intro-main}

Let $E$ be an elliptic curve over $\Q$. The Mordell--Weil theorem states that the group $E(\Q)$ of rational points of $E$ is finitely generated \cite{Mordell1922,Weil1929}; see also \cite[Theorem~VIII.6.7]{SilvermanAEC}. The rank of this group is called the rank of $E$, and we denote it by $\rk E(\Q)$. The rank is a basic invariant of $E$, but it remains poorly understood. In particular, little is known about which integers occur as the rank of infinitely many elliptic curves over $\Q$. Here we count curves up to isomorphism over $\Qbar$, so that the twists of a single curve are counted only once.

In this article, we establish that the integers five, six, and seven occur in this way.

\begin{theorem}\label{thm:main}
For each $r\in\{5,6,7\}$, there exist infinitely many elliptic curves $E$ over $\Q$, pairwise nonisomorphic over $\Qbar$, such that $\rk E(\Q)=r$.
\end{theorem}

Two elliptic curves over $\Q$ are isomorphic over $\Qbar$ if and only if they have the same $j$-invariant \cite[Proposition~III.1.4(b)]{SilvermanAEC}. Thus \cref{thm:main} states that for each $r\in\{5,6,7\}$, infinitely many distinct $j$-invariants occur among the elliptic curves over $\Q$ of rank $r$. In particular, these curves are not all twists of finitely many curves.

Each case of \cref{thm:main} comes from an explicit family. For each $r$, we fix twelve binary linear forms $L_i=a_iM+b_iN$ ($0\le i\le 11$) with integer coefficients, we set
\begin{equation*}
F=L_0L_1L_2L_3,\qquad H=L_4L_5L_6L_7,\qquad J=L_8L_9L_{10}L_{11},
\end{equation*}
and we fix nonzero integers $A$, $B$, and $C$ such that $A^2F+4BH=CJ$ identically in $M$ and $N$. For integers $M$ and $N$ with $F(M,N)H(M,N)J(M,N)\ne0$, we consider the elliptic curves
\begin{equation}\label{eq:intro-curves}
E_{M,N}\colon y^2=x^3+AFx^2-BFHx,\qquad E'_{M,N}\colon y^2=x^3-2AFx^2+CFJx,
\end{equation}
which are related by an isogeny of degree two whose kernel is generated by the point $(0,0)$. The identity $A^2F+4BH=CJ$ states that the discriminant of the quadratic polynomial $x^2+AFx-BFH$ equals $CFJ$. The discriminant of $E_{M,N}$ is then $16B^2CF^3H^2J$ (\cref{lem:datum-basic}), which is a constant times a product of powers of the twelve linear forms.

The family also comes with positive integers $c_0,\dots,c_{11}$, a modulus $W$, a residue class $(M_0,N_0)$ modulo $W$, and an open cone $\cK_r\subset\R^2$. We prove that if $(M,N)\equiv(M_0,N_0)\pmod W$, $(M,N)\in\cK_r$, and the twelve numbers $p_i=L_i(M,N)/c_i$ are distinct primes outside a fixed finite set, then $\rk E_{M,N}(\Q)=r$; see \cref{thm:r5,thm:r6,thm:r7}. Using the Green--Tao--Ziegler theorem on linear equations in primes \cite{GT2010,GT2012,GTZ2012}, we find infinitely many such pairs $(M,N)$, with infinitely many values of $M/N$. Since the $j$-invariant of $E_{M,N}$ takes each value for at most twelve values of $M/N$ (\cref{lem:j-invariant}), these curves have infinitely many $j$-invariants. The main data of the three families are listed in \cref{tab:intro-families}.

\begin{table}[ht]
\centering
\begin{tabular}{@{}crrrcrc@{}}
\toprule
$r$ & $A$ & $B$ & $C$ & $(d_E,d_{E'})$ & $W$ & Section\\
\midrule
$5$ & $1500$ & $2700$ & $57600$ & $(3,4)$ & $19984345920$ & \ref{sec:rank5}\\
$6$ & $70$ & $784$ & $1764$ & $(4,4)$ & $780373440$ & \ref{sec:rank6}\\
$7$ & $91205$ & $33124$ & $9$ & $(5,4)$ & $308403583488$ & \ref{sec:rank7}\\
\bottomrule
\end{tabular}
\caption{The three families. The integers $A$, $B$, and $C$ satisfy $A^2F+4BH=CJ$. The numbers $d_E$ and $d_{E'}$ are the dimensions of the images of the two descent maps at every specialization that we use, so that $\rk E_{M,N}(\Q)=d_E+d_{E'}-2=r$, and $W$ is the modulus of the congruence class.}
\label{tab:intro-families}
\end{table}

\section{Background}\label{sec:intro-background}

It is not known whether the ranks of elliptic curves over $\Q$ are bounded; see \cite[\S3]{PPVW2019} for the history of this question. Earlier, it was often expected that they are unbounded. In \cite{PPVW2019}, Park, Poonen, Voight, and Wood proposed a heuristic which predicts the opposite: for each $0\le r\le 20$, infinitely many elliptic curves over $\Q$ have rank $r$, and only finitely many have rank greater than $21$. The ranks in \cref{thm:main} lie in the range where this heuristic predicts infinitely many curves.

Curves of large rank are usually found among the specializations of an elliptic curve over $\Q(t)$ with many independent sections. Mestre constructed elliptic curves over $\Q(t)$ of rank at least $12$ \cite{Mestre1991}. Elkies announced constructions of higher rank in his 2007 lectures \cite{Elkies2007} and gave explicit formulas and proofs in \cite[Theorems~1--3]{Elkies2026}: an elliptic K3 surface over $\Q(t)$ of rank $17$, quadratic base changes of this surface with rank at least $18$, and a base change over a curve of genus one with rank at least $19$, which has infinitely many rational fibers of rank at least $19$. A specialization of the rank-$17$ surface is a curve over $\Q$ with at least $28$ independent rational points \cite[Theorem~2]{Elkies2026}.

For a nonisotrivial elliptic curve over $\Q(t)$, Silverman's specialization theorem \cite[Theorem~C, p.~208]{Silverman1983}, also stated in \cite[Theorem~C.20.3]{SilvermanAEC}, implies that the specialization map on sections is injective at all but finitely many rational values of $t$. Hence all but finitely many rational fibers have rank at least the rank over $\Q(t)$. Specialization yields only lower bounds. For example, Dujella and Peral proved in \cite[\S3]{DujellaPeral2020} that infinitely many curves over $\Q$ induced by rational Diophantine triples have rank at least seven, and in \cite[\S4]{DujellaPeral2024} they constructed an elliptic curve over $\Q(u)$ with torsion subgroup $\Z/4\Z$ and rank exactly six. The rank of a fiber may exceed the generic rank, and to determine it exactly one needs an upper bound as well, which in practice comes from descent. As Zywina remarks in \cite[\S1]{Zywina2026}, even with an elliptic curve of rank $18$ over $\Q(t)$, it is unclear how to find a fixed $r\ge 18$ such that infinitely many of its fibers have rank exactly $r$.

\section{Exact ranks}\label{sec:intro-exact}

For ranks zero and one, infinitely many curves over $\Q$ of exactly that rank, pairwise nonisomorphic over $\Qbar$, have been known for some time. Mazur and Rubin proved the rank-zero case over every number field \cite[Corollary~1.11 and proof of Lemma~5.4]{MazurRubin2010}. Over $\Q$, the rank-one case follows, for example, from the theorem of Bhargava and Skinner that a positive proportion of elliptic curves over $\Q$ have rank one \cite{BhargavaSkinner2014}.

An earlier use of prime values and $2$-descent to preserve a generic rank is Zywina's nonisotrivial rank-zero surface with infinitely many rank-zero fibers \cite[Theorem~1.2]{ZywinaNoJumps2025}.

In \cite{Zywina2025}, Zywina proved the corresponding statement for rank two.\footnote{Under the parity conjecture, earlier constructions of infinitely many curves over $\Q$ of rank exactly two were given by Byeon and Jeong \cite[Theorem~1.1]{ByeonJeong2016} and by Jeong \cite[Theorem~1.1]{Jeong2019}. These are twist families with constant $j$-invariant $0$ or $1728$.} His curves are $y^2=x^3-5(m+16n^2)x^2+4(m+16n^2)(m+25n^2)x$, where $m$, $m+16n^2$, and $m+25n^2$ are primes congruent to $11$ modulo $24$. He bounds the rank from above by descent through the rational $2$-isogeny, with local computations at $2$ and at the three primes, and he obtains infinitely many such triples of primes from the theorem of Tao and Ziegler on polynomial progressions in the primes \cite[Theorem~1.3]{TaoZiegler2008} (see also \cite{TaoZieglerErratum2013}). The rank-one case over every number field was proved independently by Koymans and Pagano \cite{KoymansPagano2025} and by Zywina \cite{ZywinaRankOne2025}. Savoie constructed infinitely many curves of rank two over $\Q(i)$, all with $j$-invariant $1728$ \cite{Savoie2025}.

In \cite[Theorem~1.1]{Zywina2026}, Zywina proved that for every number field $K$ and every $0\le r\le 4$, there exist infinitely many elliptic curves over $K$ of rank $r$, pairwise nonisomorphic over $\overline K$. For each $r$, he uses an explicit nonisotrivial elliptic curve of rank $r$ over $K(T)$ with a point of order two, and he shows that infinitely many of its specializations have rank $r$; we describe his method in \cref{sec:intro-method}. Our proof of \cref{thm:main} follows his strategy. Over $\Q$, \cref{thm:main} extends the range $0\le r\le 4$ of \cite[Theorem~1.1]{Zywina2026} to $r\in\{5,6,7\}$. We do not consider other number fields.

Two recent results are close to \cref{thm:main}. In \cite[Theorem~1.1]{Savoie2026}, Savoie proved that for each $r\in\{4,6\}$, there exist infinitely many elliptic curves over $\Q(i)$ of rank exactly $r$, pairwise nonisomorphic over $\Q(i)$, not obtained by base change from $\Q$, and all with $j$-invariant $1728$. His curves are quartic twists $y^2=x^3+bx$. Their points come from factorizations $b=ef$ with $e+f$ a square, the upper bounds come from descent through the endomorphism $1+i$, and Kai's theorem \cite{Kai} on prime values of linear forms over number fields supplies the specializations. Since the $j$-invariant is constant, all of these curves are isomorphic over $\Qbar$. Compared with \cref{thm:main}, the ground field and the notion of infinitude differ: our curves are defined over $\Q$ and have infinitely many $j$-invariants. Both works use explicit points whose classes have prescribed prime supports, and both express the local conditions in the descent through quadratic residue symbols. Savoie uses the graph and matrix descriptions of Kling and Savoie \cite[Theorem~5.2 of the cited preprint]{KlingSavoie2024}. In his rank-six family, he reduces the pairwise symbols to one undetermined sign and proves that both resulting Laplacians have the same kernel \cite[Lemmas~4.1--4.2]{Savoie2026}. In \cite[Proposition~6.1]{Lu2026}, Lu constructed five independent sections over the function field of a curve of genus one with infinitely many rational points, and obtained infinitely many elliptic curves over $\Q$, pairwise nonisomorphic over $\Qbar$, of rank at least five. As Lu notes, these results give lower bounds and do not determine the ranks.

\section{The method and its sources}\label{sec:intro-method}

We now describe the ingredients of the proof of \cref{thm:main} and their origins. \Cref{sec:overview} gives a more detailed overview.

\paragraph{Descent by a rational $2$-isogeny.} For a pair $(M,N)$ with $FHJ\ne0$, let $\phi\colon E_{M,N}\to E'_{M,N}$ be the isogeny between the curves in \eqref{eq:intro-curves}, and let $\dual$ be its dual. The descent map $\delta_E$ sends the identity of $E_{M,N}(\Q)$ to the trivial class in $\sqgp{\Q}$, the point $(0,0)$ to the class of $-BFH$, and every other point $(x,y)$ to the class of $x$. It is a homomorphism with kernel $\dual(E'_{M,N}(\Q))$. The map $\delta_{E'}$ on $E'_{M,N}(\Q)$ is defined in the same way and has kernel $\phi(E_{M,N}(\Q))$. If $d_E$ and $d_{E'}$ denote the dimensions over $\F_2$ of the images of $\delta_E$ and $\delta_{E'}$, then $\rk E_{M,N}(\Q)=d_E+d_{E'}-2$. Descent is the method by which Mordell and Weil proved finite generation \cite{Mordell1922,Weil1929}. The form of it that we use, for a rational $2$-isogeny, is classical. We follow Tate's 1961 Haverford lectures as expanded in \cite[\S\S3.4--3.6]{SilvermanTate2015}, together with \cite[Example~X.4.8 and Proposition~X.4.9]{SilvermanAEC} and \cite[\S14]{Cassels1991}. The rank formula in the form above is \cite[Lemma~3.1]{Zywina2026}; see also \cite[\S3.6]{SilvermanTate2015}.

We bound $d_E$ and $d_{E'}$ from below by the classes of explicit rational points. For the upper bounds, we use that for every place $v$ of $\Q$, the image of $\delta_E(E_{M,N}(\Q))$ in $\sqgp{\Q_v}$ lies in the image of the local descent map on $E_{M,N}(\Q_v)$, and similarly for $E'_{M,N}$. These are the local conditions that define the Selmer groups of $\phi$ and $\dual$ \cite[\S X.4]{SilvermanAEC}. We impose only the conditions that we need, at finitely many places, and show that the resulting upper bounds agree with the lower bounds.

\paragraph{Local conditions as linear equations.} At the specializations we use, every class in the image of $\delta_E$ is represented by a product of $-1$, a few fixed primes, and some of the primes $p_i$ that divide $BFH$ (\cref{lem:support}). Once we fix these generators, we may regard a class as a vector over $\F_2$, and the local condition at a prime $p_i$ is a linear equation in this vector whose coefficients are Legendre symbols modulo $p_i$. The local conditions at the fixed places are linear equations as well. Such matrices occur in R\'edei's arithmetic treatment \cite[\S5 and p.~59, footnote~11]{Redei1934} of the R\'edei--Reichardt formula for the $4$-rank of the narrow class group of a quadratic field \cite{RedeiReichardt1934}. For elliptic curves, Monsky gave a binary-matrix description of full $2$-descent for the congruent-number curves in his appendix to \cite[pp.~365--369, especially p.~366]{HeathBrownMonsky1994}. Feng used quadratic-residue graphs and their binary Laplacians in descent by a $2$-isogeny \cite[\S\S2--3]{Feng1996}. Xiong and Zaharescu treated the Selmer groups of a $2$-isogeny of these curves and of its dual in a similar way \cite[\S3.1, Theorem~3]{XiongZaharescu2009}, and Kane and Klagsbrun constructed the corresponding matrices for a rational $2$-isogeny and its dual in families of quadratic twists \cite[\S4]{KaneKlagsbrun2017}. In these congruent-number and quadratic-twist treatments, the primes are the prime factors of a twisting parameter, and the Legendre symbols among them vary with the parameter. In our families, a congruence condition on $(M,N)$ determines every entry of the matrices as a function of a single undetermined sign, which we discuss below.

\paragraph{Prime values and descent.} To control the primes of bad reduction of $E_{M,N}$, we specialize at pairs $(M,N)$ where the twelve linear forms take prime values, up to the fixed factors $c_i$. Prime values of polynomials have long been used for this purpose in the fibration method. Colliot-Th\'el\`ene and Sansuc used Schinzel's Hypothesis~H \cite{SchinzelSierpinski1958}, which is a conjecture, to study the Hasse principle and weak approximation on fibrations \cite{ColliotTheleneSansuc1982}; see also \cite{ColliotTheleneSwinnertonDyer1994}. Swinnerton-Dyer \cite{SwinnertonDyer1995} introduced this approach for certain genus-one pencils over $\Q$ whose Jacobians have full rational $2$-torsion. Colliot-Th\'el\`ene, Skorobogatov, and Swinnerton-Dyer \cite[Introduction]{CSS1998GenusOne} developed it further for suitable pencils over number fields, choosing primes so as to control the $2$-Selmer groups. These results assume Hypothesis~H and the finiteness of the relevant Tate--Shafarevich groups. Bender and Swinnerton-Dyer \cite[Theorems~1--2]{BenderSwinnertonDyer2001} extended this method to certain pencils of $2$-coverings whose generic Jacobians have exactly one nonzero rational $2$-torsion point. Colliot-Th\'el\`ene \cite[Theorems~A--B and \S1.4]{ColliotThelene2001} recast their solubility criteria in terms of the vertical Brauer--Manin obstruction, using the Selmer groups of a $2$-isogeny and its dual; these results retain the same conjectural assumptions. Skorobogatov and Swinnerton-Dyer \cite{SkorobogatovSwinnertonDyer2005} developed a version of this descent for certain Kummer surfaces which still assumes the finiteness of the relevant Tate--Shafarevich groups but does not require Hypothesis~H. When the polynomials are linear forms in two variables, as they are for a fibration over the projective line whose degenerate fibers lie over rational points, the Green--Tao--Ziegler theorem can take the place of Hypothesis~H. Browning, Matthiesen, and Skorobogatov used additive combinatorics together with descent over $\Q$ to treat pencils of conics and quadrics with rational degenerate fibers \cite[Theorems~1.1--1.2]{BrowningMatthiesenSkorobogatov2014}. Harpaz, Skorobogatov, and Wittenberg adapted the fibration method over $\Q$ using simultaneous linear prime values in \cite{HSW2014}, for suitable fibrations of this kind. Kai \cite{Kai} extended the Green--Tao--Ziegler theorem to prime elements of number fields. Koymans and Pagano combine four prime-valued linear forms, prescribed quadratic characters at fixed auxiliary primes, and Hilbert reciprocity to control a twist's Selmer group \cite[\S3, Theorems~3.3 and~3.5]{KoymansPagano2025}.

\paragraph{Linear equations in primes.} Green and Tao \cite{GT2010} proved an asymptotic formula for the number of integer points in a convex region at which a system of affine-linear forms of finite complexity takes prime values, assuming two conjectures: the M\"obius--nilsequences conjecture and the inverse conjecture for the Gowers norms. The first was proved by Green and Tao \cite[Theorem~1.1]{GT2012}, and the cases of the second that were not previously known were proved by Green, Tao, and Ziegler \cite[Theorem~1.3]{GTZ2012}, whose later corrections are recorded in the bibliography. We use this theorem for twelve affine-linear forms in two variables with pairwise nonproportional linear parts. Such a system has complexity at most ten by \cite[Lemma~1.6]{GT2010}, and \cref{sec:primes} states the form of the theorem that we need. It produces prime values at points of a prescribed congruence class and cone, provided that there is no local obstruction, but it gives no control of the Legendre symbols among the resulting primes.

\paragraph{Zywina's strategy.} In \cite{Zywina2026}, Zywina starts from a nonisotrivial elliptic curve $y^2=x^3+ax^2+bx$ over $K(T)$ whose discriminant is a product of linear factors over $K$, so that its bad fibers lie over finitely many points $e\in K$. Sections of this curve and of the $2$-isogenous curve over $K(T)$ bound the images of the descent maps from below at all but finitely many specializations. He chooses a finite set $S$ of primes of $K$ containing the denominators of the bad parameters, together with two auxiliary primes $\mathfrak q_1,\mathfrak q_2$ outside $S$. Kai's theorem \cite[Proposition~13.2]{Kai} then provides $S$-integers $m$ and $n$ such that the elements $m-en$ are pairwise nonassociate prime elements of the ring of $S'$-integers, where $S'=S\cup\{\mathfrak q_1,\mathfrak q_2\}$, and $(m,n)$ satisfies prescribed local conditions; he specializes at $t=m/n$ \cite[Proposition~5.9]{Zywina2026}. Using local computations at these primes, auxiliary primes, global reciprocity, and the Selmer ratio of the isogeny, Zywina shows that the Selmer groups of the specialized isogenies are spanned by the classes of the sections. The rank of the specialization then equals the rank over $K(T)$.

The axiomatic form of this argument is \cite[Theorem~4.1]{Zywina2026}. We write $\mathcal A$ for the set of additive fibers, and $\mathcal M$ (respectively $\mathcal M'$) for the set of multiplicative fibers of type $\mathrm I_{2n}$ on the first curve and $\mathrm I_n$ on the second (respectively the reverse). Besides conditions on the bad fibers, Theorem~4.1 of \cite{Zywina2026} requires that the images of the two descent maps over $K(T)$ have dimensions at least $|\mathcal A|+|\mathcal M|-1$ and $|\mathcal A|+|\mathcal M'|-1$. These conditions force the rank over $K(T)$ to equal $2|\mathcal A|+|\mathcal M|+|\mathcal M'|-4$, which is the upper bound $\deg\mathcal N-4$ of Shioda \cite[Corollary~2]{Shioda1992}, where $\mathcal N$ is the conductor of the elliptic surface. In characteristic zero, Shioda's stronger bound $\rk E(\overline K(T))\le\deg\mathcal N-4-2p_g$ \cite[Corollary~3, equation~(1.13')]{Shioda1992} shows that equality in $\deg\mathcal N-4$ forces the geometric genus $p_g$ to be zero; see also \cite[\S2]{Zywina2026}. Zywina's surfaces are rational elliptic surfaces with a point of order two, whose rank is at most four by \cite[Corollary~2.1]{OguisoShioda1991}. For $r=4$, he uses the pair
\begin{equation}\label{eq:intro-zywina}
\begin{aligned}
E_Z\colon y^2&=x^3-70(T^2-25^2)x^2+2^4 7^2(T^2-11^2)(T^2-25^2)x,\\
E'_Z\colon y^2&=x^3+140(T^2-25^2)x^2+2^2 3^2 7^2(T^2-25^2)(T^2-39^2)x
\end{aligned}
\end{equation}
of \cite[\S6.5]{Zywina2026}, whose six bad fibers lie over $T=\pm11,\pm25,\pm39$.

\paragraph{Conic base change.} Our families arise by base change of degree two. For the rank-five family, we start from the curve $E'_Z$ in \eqref{eq:intro-zywina}, and for the rank-six family from $E_Z$. For the rank-seven family, we start from a related $2$-isogenous pair of the same form, with bad fibers over $T=\pm1,\pm v,\pm w$ for certain rational numbers $v$ and $w$. In each case, the base change is along a conic over which an additional point becomes rational; \cref{sec:overview-origin} explains how it produces the twelve forms $L_i$ and the identity $A^2F+4BH=CJ$, and \cref{app:origins} gives the details. Making a quadratic section rational by a base change of degree two is a classical way to raise the rank of an elliptic surface; see Mestre's conic construction \cite[\S2.2]{Mestre1991} and Elkies's discussion \cite[Lecture~III, p.~2003]{Elkies2007}. Dujella, Kazalicki, and Peral describe the corresponding divisor and square conditions for models $y^2=x^3+ax^2+bx$ in \cite[\S1.1]{DujellaKazalickiPeral2021}. Salgado gives general results on new independent sections from moving multisections \cite[Corollary~4.3]{Salgado2012}, and Loughran and Salgado develop this method using conic bundles \cite[\S1.2 and Theorem~2.4]{LoughranSalgado2022}. In \cite[\S4]{DujellaPeral2024}, Dujella and Peral raise the rank of an elliptic curve over $\Q(u)$ from four to six by imposing two conditions, each of which defines a conic. Such a base change need not raise the rank. For the rank-seven family, the point supplied directly by the conic has a square $x$-coordinate, so its class under $\delta_E$ is trivial, and the additional class comes from another point (\cref{app:origins}). For this reason, the proofs do not use the construction. We verify the new classes directly at every specialization, and \cref{app:origins} only explains where the families come from.

\paragraph{What is new.} The proof of \cref{thm:main} differs from the maximal-rank argument of \cite{Zywina2026} in the following ways.
\begin{itemize}
\item \emph{Three families.} These are the families of \cref{tab:intro-families}, with their points, congruence classes, and cones.
\item \emph{Explicit matrices instead of the criterion of maximal rank.} After the base change, each family has twelve bad fibers: four additive fibers at the zeros of $F$, and eight multiplicative fibers at the zeros of $H$ and $J$, which form the sets $\mathcal M$ and $\mathcal M'$ respectively. The number $2|\mathcal A|+|\mathcal M|+|\mathcal M'|-4$ is then $12$, while the ranks we prove are $5$, $6$, and $7$. Since infinitely many fibers of each family have rank $r\le 7$, Silverman's theorem \cite[Theorem~C, p.~208]{Silverman1983} implies that the rank of the family over $\Q(M/N)$ is at most $r$. Thus the families do not have maximal rank, and \cite[Theorem~4.1]{Zywina2026} does not apply. (The fiber types and K3 classification of these elliptic surfaces are proved in \cref{rem:elliptic-surfaces}, using \cite{Tate1975,SchuttShioda2010}.) Instead, for a fixed congruence class and cone, we write the local conditions at $2$, at $3$, at one or two further fixed primes, at the real place, and at the twelve primes $p_i$ as the rows of explicit matrices over $\F_2$ with twelve or thirteen columns. Their ranks show that the images of $\delta_E$ and $\delta_{E'}$ are spanned by the classes of explicit points. We display the matrices with labeled rows and columns, explain each entry through tables of residues, and give rank certificates that can be checked by hand.
\item \emph{One free sign.} The congruence class determines every entry of the matrices as a function of the sign $\eps=\leg{N}{p_0}$, which it does not determine. When $\eps$ changes sign, every row coming from one of the primes $p_i$ changes by the same vector, and both values of $\eps$ give the same upper bound (\cref{lem:eta-rank}). The analytic input is therefore used only to produce prime values in a congruence class and a cone, and we need no condition on quadratic residue symbols among the growing primes. The related one-sign reduction and common-kernel mechanism over $\Q(i)$ in \cite[Lemmas~4.1--4.2]{Savoie2026} is a precedent for this step.
\item \emph{The ground field.} The proof takes place over $\Q$ and uses the Green--Tao--Ziegler theorem directly, instead of Kai's extension to number fields.
\end{itemize}
The proof of \cref{thm:main} is unconditional. It does not use the finiteness of Tate--Shafarevich groups, the parity conjecture, analytic ranks, or a specialization theorem: both the upper and the lower bounds are proved directly at each prime specialization.

\section{Computer verification}\label{sec:intro-verification}

In \cref{part:families}, we use finite arithmetic of the following kinds:
\begin{itemize}
\item the identity $A^2F+4BH=CJ$, and the identities stating that the listed points lie on the curves, which are homogeneous polynomial identities of degree at most twelve in $(M,N)$;
\item the values of $W$, $M_0$, and $N_0$ given by the Chinese remainder theorem;
\item the factorizations of the coefficients $b_i$ and of the determinants $a_ib_j-a_jb_i$;
\item the residues of the primes $p_i$ modulo small primes, and the Legendre symbols that they determine;
\item the ranks of matrices over $\F_2$ with twelve or thirteen columns.
\end{itemize}
We state each of these claims precisely, so that a reader can check it by hand or with any computer algebra system. For instance, a homogeneous identity of degree at most twelve in $(M,N)$ holds if it holds at $N=1$ for thirteen distinct values of $M$. For the ranks, we display the matrices and give certificates in the form of triangular minors. \Cref{app:verification} describes the supplied scripts, which use exact arithmetic, and lists the claims that they check. The uniform arguments, which apply to every pair $(M,N)$ in an infinite progression, and the analytic input are proved in the text and are not checked by computer.

\section{Organization and notation}\label{sec:intro-organization}

\Cref{sec:overview} gives an overview of the proof, including a dimension count for one of the families. \Cref{part:tools} develops the general tools. \Cref{sec:descent} recalls descent by a rational $2$-isogeny over $\Q$ and over its completions, proves the rank formula (\cref{lem:rank-formula}) and the local lemmas that we need, and shows how local conditions become rows of a matrix over $\F_2$ (\cref{lem:upper-bound}). \Cref{sec:families} introduces twelve-form data (\cref{def:datum}) and their prime specializations, computes the local images at the twelve primes $p_i$ and the Legendre symbols among them, and proves the rank criterion of \cref{prop:rank-criterion}. \Cref{sec:primes} states the Green--Tao--Ziegler theorem in the form that we use and proves \cref{prop:infinitely-many}, which gives infinitely many prime specializations in a congruence class and a cone. \Cref{part:families} treats the three families. \Cref{sec:rank5,sec:rank6,sec:rank7} prove \cref{thm:r5,thm:r6,thm:r7} by checking the hypotheses of \cref{prop:rank-criterion,prop:infinitely-many}, and \cref{sec:proof} deduces \cref{thm:main}. \Cref{app:verification} describes the computer verification, and \cref{app:origins} explains how the three families were found.

Throughout, for a field $K$ of characteristic zero, we write $\sqgp{K}$ for the group of squareclasses of $K$, which we regard as a vector space over $\F_2$, and $\sq{u}$ for the class of $u\in K^\times$. For a prime $p$, we write $v_p$ for the $p$-adic valuation on $\Q$ and on $\Q_p$, and we set $\Q_\infty=\R$. For an odd prime $p$ and an integer $z$ prime to $p$, we write $\leg zp$ for the Legendre symbol. At a prime specialization, for $0\le i\le 11$ and an integer $z$ prime to $p_i$, we define $\lambda_i(z)\in\F_2$ by $\leg{z}{p_i}=(-1)^{\lambda_i(z)}$, so that $\lambda_i(zz')=\lambda_i(z)+\lambda_i(z')$. We write the free sign as $\eps=\leg{N}{p_0}=(-1)^\eta$ with $\eta\in\F_2$. The remaining notation, including the curves $E_{\alpha,\beta}$, the isogeny $\phi$ and its dual $\dual$, and the descent maps, is introduced in \cref{sec:descent,sec:families}.
